
% ====================
\section{Day 1}
% ====================

\subsection{Readings}

\begin{itemize}
    \item Albert Einstein and Leopold Infeld,
    \textit{The Evolution of Physics}.
    \item Charles Darwin's notebooks.
    \item Isaac Newton's ``Waste Book''.
    \item John Gerhart and Marc Kirschner,
    \textit{Cells, Embryos, and Evolution}.
\end{itemize}


% ====================
\subsection{The scale of things}
% ====================

When a quantity is known only to lie somewhere between two scales,
$x_{\min}$ and $x_{\max}$, a useful order-of-magnitude estimate is the geometric mean,

\[
    x_{\mathrm{est}} \sim \sqrt{x_{\min}x_{\max}}.
\]

This is especially natural when the uncertainty spans orders of magnitude.
On a logarithmic scale,

\[
    \log x_{\mathrm{est}} = \frac{1}{2} \left( \log x_{\min} + \log x_{\max} \right),
\]

so the estimate lies halfway between the two bounds in log-space.

Various biological examples with full breakdowns are available at \url{https://book.bionumbers.org}.


% ====================
\subsection{Gorter--Grendel experiment and the lipid bilayer}
% ====================

Gorter and Grendel used red blood cells to infer that cellular membranes
consist of a lipid bilayer.

The basic logic was:

\begin{enumerate}
    \item Take a known number of red blood cells.
    \item Estimate the total surface area of those cells.
    \item Extract the lipids from the cells.
    \item Spread the extracted lipids into a monolayer.
    \item Measure the area occupied by this lipid monolayer.
\end{enumerate}

The lipid monolayer occupied approximately twice the area of the cell
surfaces. This suggests that the membrane contains two layers of lipid:

\[
    A_{\mathrm{lipid}} \approx 2 A_{\mathrm{cell}}.
\]

Hence the membrane is approximately a lipid bilayer.


% ====================
\subsection{The cell as a bag of proteins}
% ====================

A useful coarse-grained picture of the cell is as a system whose biomass
is dominated by proteins and whose growth is therefore closely related
to protein production.

An order-of-magnitude estimate discussed in the lecture is

\[
    \frac{\text{number of protein molecules}}{\text{number of mRNA molecules}} \sim 10^3.
\]

Ribosomes synthesize proteins, including the ribosomal proteins required to construct new ribosomes. 
Thus the translational machinery participates in its own production. 
At a coarse-grained level this creates an \emph{autocatalytic} process: the machinery responsible for growth produces more of the machinery responsible for growth.

This provides an intuitive route to exponential cellular growth.

% ====================
\subsubsection*{Simple growth law}
% ====================

Let $M(t)$ denote the cell mass at time $t$.

Assume that the rate at which new mass is produced is proportional to the mass already present. 
Over a short time interval $\Dt$, 
\[
    M(t+\Dt) = M(t) + kM(t)\Dt,
\] 
where $k$ is the fractional growth rate. 
Rearranging, 
\[
    \frac{M(t+\Dt)-M(t)}{\Dt} = kM(t).
\] 
Taking the limit $\Dt\rightarrow 0$ gives 
\[
    \frac{\dd M}{\dd t} = kM.
\] 
The solution is 
\[ 
    M(t) = M(0)e^{kt}.
\] 
Thus a constant fractional growth rate produces exponential growth.


\subsubsection*{Assumptions}

This simple model assumes that

\begin{itemize}
    \item the cell can be described by a single coarse-grained variable $M(t)$;
    \item the amount of active protein-production machinery is
    approximately proportional to cell mass;
    \item the protein-production rate per unit mass is approximately
    constant;
    \item nutrients and other resources are not limiting over the time interval considered;
    \item degradation and other losses are either negligible or already incorporated into the effective rate constant $k$;
    \item fluctuations in individual molecular events can be neglected, so that a deterministic continuum description is appropriate.
\end{itemize}


% ====================
\subsection{Biological turnover}
% ====================

Biological systems continually replace their components. Examples include
\begin{itemize}
    \item red blood cells,
    \item photoreceptor components,
    \item proteins and RNA molecules,
    \item cells in renewing tissues.
\end{itemize}

A simple birth-death model illustrates how turnover can generate a steady state even when individual molecules are continuously produced and destroyed.


% ====================
\subsection{mRNA production and degradation}
% ====================

Let $m(t)$ denote the number of mRNA molecules in a cell.

Suppose that
\begin{itemize}
    \item mRNA molecules are produced at a constant rate $r$;
    \item each existing mRNA has a constant probability per unit time $\gamma$ of being degraded.
\end{itemize}

For example, a typical order-of-magnitude transcription rate might be
\[
    r \sim 1 \text{ mRNA per few minutes},
\]
while the characteristic mRNA lifetime is
\[
    \tau = \frac{1}{\gamma}.
\]
Over a small time interval $\Dt$,
\begin{align*}
    m(t+\Dt) &= m(t) + \text{mRNAs produced} - \text{mRNAs degraded}
    \\
    &= m(t) + r\Dt - \gamma m(t)\Dt.
\end{align*}
Therefore,
\[
    \frac{m(t+\Dt)-m(t)}{\Dt} = r-\gamma m(t).
\]
Taking $\Dt \rightarrow 0$,
\[
    \frac{\dd m}{\dd t} = r-\gamma m.
\]

At steady state,
\[
    \frac{\dd m}{\dd t}=0,
\]
so
\[
    m_{\mathrm{ss}} = \frac{r}{\gamma}.
\]
The full solution is

\[
    m(t) = \frac{r}{\gamma} + \left( m(0)-\frac{r}{\gamma} \right) e^{-\gamma t}.
\]
Thus the system approaches its steady-state mRNA number on the characteristic timescale
\[
    \tau = \frac{1}{\gamma}.
\]

% ====================
\subsubsection*{Assumptions}
% ====================

The deterministic model assumes that

\begin{itemize}
    \item transcription occurs at a constant average rate $r$;
    \item degradation is first-order: each mRNA has the same constant degradation probability per unit time;
    \item mRNA molecules act independently;
    \item $r$ and $\gamma$ do not vary with time;
    \item discrete molecular fluctuations are ignored.
\end{itemize}

At the level of individual molecules, transcription and degradation are
stochastic processes. The equation above describes the mean behavior.

